\clearpage\thispagestyle{empty}
\pdfbookmark[0]{Quellen zur Besprechung von Rossels zweitem Band}{rossel1809-sources}
\begin{center}\textbf{Quellen zur Besprechung von Rossels zweitem Band}\end{center}
\noindent Diese KI-übersetzte Fassung in modernem Deutsch folgt Carl Friedrich Gauß, \href{https://gdz.sub.uni-goettingen.de/id/PPN236007467}{\emph{Werke}, Band~VI} (Göttingen, 1874), S.~520–522. Die Besprechung erschien zuerst ohne Verfassernamen unter Paris in den \emph{Göttingischen gelehrten Anzeigen}, Stück~167, 21.~Oktober 1809, \href{https://images.p.sub.uni-goettingen.de/iiif/image/gdz:PPN319721507_1809:00001661/full/1650,/0/default.jpg}{S.~1657}–\href{https://images.p.sub.uni-goettingen.de/iiif/image/gdz:PPN319721507_1809:00001667/full/1650,/0/default.jpg}{1663}. Sie behandelt den zweiten Band über die Reise von d’Entrecasteaux, der auf die Suche nach La Pérouse geschickt wurde. Der französische Titel nennt die Veröffentlichung auf Anordnung des Kaisers und Königs und die Bearbeitung durch M. de Rossel, einen früheren Kapitän zur See. Der \href{https://catalogue.bnf.fr/ark:/12148/cb43383186m}{Katalog der Bibliothèque nationale de France} führt Rossel als wissenschaftlichen Herausgeber.
\par\smallskip\noindent In der Werkausgabe steht ein Punkt nach \emph{ausmacht} und ein apostrophähnliches Zeichen vor \emph{den}. Im Erstdruck steht dort ein Komma, und das Zeichen fehlt. Bei den Schwingungszeiten stehen im Erstdruck Doppelstriche, in der Werkausgabe ein hochgestelltes \emph{s}. Die Übertragung folgt der Werkausgabe.
\par\smallskip\noindent In Abschnitt V zitiert Gauß Rossel auf Französisch.
\begin{quote}
Tout ce qui peut faciliter une observation aussi délicate, est précieux; nous n'en connoissons pas qui exige plus de dispositions naturelles: aussi nous sommes bien convaincus que, lorsque les savans auront fait, pour la perfection de cette méthode, tout ce que leur génie pourra leur inspirer, le talent de l'observateur aura encore beaucoup à y ajouter.
\end{quote}
\noindent Übersetzung und Satz mit KI erstellt.
\clearpage
